% A4 review layout: use the page instead of the class's small-format text block.
% Keep review line spacing; limit whitespace around figures and the title.
\geometry{reset,a4paper,left=22mm,right=24mm,top=20mm,bottom=22mm,
          bindingoffset=0mm,footskip=10mm}
\raggedbottom
\setlength{\floatsep}{10pt plus 2pt minus 2pt}
\setlength{\textfloatsep}{10pt plus 2pt minus 2pt}
\setlength{\intextsep}{10pt plus 2pt minus 2pt}
\renewcommand{\topfraction}{.95}
\renewcommand{\bottomfraction}{.90}
\renewcommand{\textfraction}{.05}
\renewcommand{\floatpagefraction}{.80}
\makeatletter
% Float-only pages start at the top instead of being vertically centred.
\setlength{\@fptop}{0pt}
\setlength{\@fpsep}{10pt plus 2pt minus 2pt}
\usepackage{etoolbox}
\patchcmd{\@maketitle}{\vskip21pt}{\vskip0pt}{}{}
% The class also inserts an empty opening box and an unused article-type row.
\patchcmd{\@maketitle}{\newpage\null}{\newpage}{}{\PackageError{synkit-layout}{Cannot remove title opening box}{Check the class title definition.}}
\patchcmd{\@maketitle}{{\hbox to \textwidth{{\Artcatfont\ArtType\hfill}\par}}}{}{}{\PackageError{synkit-layout}{Cannot remove unused article-type row}{Check the class title definition.}}
\patchcmd{\@maketitle}{\vskip20pt}{\vskip0pt}{}{\PackageError{synkit-layout}{Cannot remove title padding}{Check the class title definition.}}
\pretocmd{\Titlefont}{\def\baselinestretch{1}}{}{}
\patchcmd{\@maketitle}{\vskip36pt}{\vskip12pt}{}{}
\makeatother

% Capability marks for the software-scope table.  Defined here, in the shared
% preamble, because xr-hyper imports main.aux (and its caption) into the SI.
% Shape carries the meaning; the marks are distinguishable without colour.
\newcommand{\yes}{\ensuremath{\checkmark}}
\newcommand{\no}{\ensuremath{\times}}
\newcommand{\pmark}{\ensuremath{(\checkmark)}}
\newcommand{\rot}[1]{\rotatebox[origin=l]{90}{#1}}

% Shared TikZ figure styling, loaded after the document packages.
\AtEndPreamble{%
% Load in the document preamble. XeLaTeX/LuaLaTeX and an installed Arial
% family are required. Liberation Sans, metric-compatible with Arial, is the
% fallback where Arial is not installed.
\RequirePackage[no-math]{fontspec}
% Figures now use the text face, Latin Modern Roman, in a plain black-and-white
% print style. The family keeps its old name, \scArial, so the figure sources
% need no change. Only the visual abstract (a PDF from visual_abstract.py)
% keeps Arial and colour.
\newfontfamily\scArial{Latin Modern Roman}[NFSSFamily=scArial,Ligatures=TeX]
% Keep manuscript mathematics unchanged; select Arial letters and numerals
% only inside figures. TeX supplies specialist mathematical symbols.
\RequirePackage[subdued,italic]{mathastext}
\MTDeclareVersion{synarial}{TU}{scArial}{m}{n}
\SetMathAlphabet{\Mathsf}{synarial}{TU}{scArial}{m}{n}
\SetMathAlphabet{\Mathtt}{synarial}{TU}{scArial}{m}{n}

\usetikzlibrary{positioning,calc,fit,backgrounds,arrows.meta,shapes.geometric,
                shapes.misc,decorations.pathreplacing,patterns,math}

% --- canvas ---------------------------------------------------------------
% 372 pt text width minus the 1 pt standalone border on each side.
\newlength{\scW}\setlength{\scW}{13.0326cm}
\providecommand{\scWnum}{13.0326}
% \scCanvas{<height>} pins the picture to the full text width, y running down
% from 0, so every figure shares one coordinate convention.  It *sets* the
% bounding box rather than contributing to it, so the emitted page is exactly
% the text width and every figure is included at 1:1 -- otherwise a stray
% invisible extent (pgfplots reserves a little around an axis) rescales one
% figure relative to the others and the same \scSmall prints at two sizes.
% Anything drawn outside the canvas is therefore cropped: keep content inside,
% and check with
%     \providecommand{\scCanvas}[1]{\path[draw=red] (0,0) rectangle (\scWnum,-#1);}
\providecommand{\scCanvas}[1]{%
  \path[use as bounding box] (0,0) rectangle (\scWnum,-#1);}

% --- palette --------------------------------------------------------------
% Canonical companion-paper values (../../../../Style/style.tex).
% Black and white: the former hues map to ink and grays, so every figure
% prints the same in grayscale. Shape, fill, and dashing carry the meaning.
\definecolor{npgOrange}{gray}{0.35}
\definecolor{npgSky}{gray}{0.55}
\definecolor{npgGreen}{gray}{0}
\definecolor{npgYellow}{gray}{0.75}
\definecolor{npgBlue}{gray}{0}
\definecolor{npgVermillion}{gray}{0.45}
\definecolor{npgPurple}{gray}{0.4}
\definecolor{figInk}{gray}{0}
\definecolor{figMuted}{gray}{0.3}
\definecolor{figHair}{gray}{0.72}
\definecolor{figWash}{gray}{1}

\colorlet{scInk}{figInk}              % primary text
\colorlet{scInkM}{figMuted}           % secondary text
\colorlet{scPaper}{figWash}           % quiet cards and note strips
\colorlet{scSpF}{npgSky!12}           % species fill
\colorlet{scSpE}{npgBlue}             % species edge
\colorlet{scRxF}{npgOrange!12}        % reaction fill
\colorlet{scRxE}{npgOrange!88!black}  % reaction edge
\colorlet{scReac}{npgBlue!78!black}   % reactant edge
\colorlet{scProd}{npgGreen!82!black}  % product edge
\colorlet{scHiF}{npgOrange!13}        % highlight fill
\colorlet{scHiE}{npgOrange!82!black}  % highlight edge
% Stage / accent tints.  stDE remains available for the fifth plot series.
\colorlet{stAF}{npgGreen!9}\colorlet{stAE}{npgGreen}
\colorlet{stBF}{npgOrange!11}\colorlet{stBE}{npgOrange}
\colorlet{stCF}{npgBlue!9}\colorlet{stCE}{npgBlue}
\colorlet{stDF}{npgPurple!10}\colorlet{stDE}{npgPurple!85!black}
\colorlet{stEF}{npgVermillion!9}\colorlet{stEE}{npgVermillion}

% Semantic aliases.  Figures should speak in these, not in stAE/stBE/...
\colorlet{scAccent}{stAE}\colorlet{scAccentF}{stAF}   % structure / normal
\colorlet{scMark}{stBE}\colorlet{scMarkF}{stBF}       % highlight / certificate
\colorlet{scRef}{stCE}\colorlet{scRefF}{stCF}         % secondary series
\colorlet{scAlert}{stEE}\colorlet{scAlertF}{stEF}     % failure / violation
\colorlet{scRule}{figHair}                            % card and axis hairlines
\colorlet{scRuleL}{figHair!55}                        % grid, dividers

% --- text styles ----------------------------------------------------------
% Three sizes, fixed in points so they do not drift with the document class.
% Arial is shared by labels, headings, panel badges, and chemistry labels.
\providecommand{\scTitle}{\scArial\fontsize{8}{9.6}\selectfont\color{scInk}}
\providecommand{\scBody}{\scArial\fontsize{7}{8.4}\selectfont\color{scInk}}
\providecommand{\scSmall}{\scArial\fontsize{6.5}{7.8}\selectfont\color{scInk}}
\providecommand{\scSub}[1]{{\scSmall\itshape #1}}
\providecommand{\scPanel}[1]{{\scTitle(#1)}}
% Numerals that carry a result: ink, not colour, so the eye reads the value.
\providecommand{\scNum}{\scArial\fontsize{10}{11}\selectfont\bfseries\color{scInk}}

% --- shared TikZ styles ---------------------------------------------------
\tikzset{
  % text never hyphenates inside a figure
  every node/.append style={execute at begin node={%
      \hyphenpenalty=10000 \exhyphenpenalty=10000 \hbadness=10000 }},
  syn/.style={font=\scArial,color=scInk,
              execute at begin picture={\MTversion*{synarial}},
              >={Stealth[round,length=2.4mm,width=1.9mm]}},
  % containers ------------------------------------------------------------
  % One card chrome for the whole paper: white, 3 pt radius, one hairline.
  % Plain print style: panels, cards, notes and chips draw no frame.
  panel/.style={fill=none,draw=none},
  card/.style={fill=none,draw=none,inner sep=5pt,align=left,anchor=north west},
  note/.style={fill=none,draw=none,inner xsep=6pt,inner ysep=3.5pt,align=left,
               font=\scSmall},
  chip/.style={fill=none,draw=none,inner xsep=4pt,inner ysep=2pt,font=\scSmall},
  % Circular letter badge: the same accessible panel marker used in Style/*.
  hbullet/.style={circle,fill=scInk,text=white,draw=none,
                  font=\scArial\fontsize{6}{6}\selectfont\bfseries,
                  minimum size=4.2mm,inner sep=0pt},
  ctag/.style={font=\scTitle,anchor=north west,text=scInk},
  % CRN primitives --------------------------------------------------------
  species/.style={rounded rectangle,fill=scSpF,draw=scSpE,line width=0.45pt,
                  inner xsep=2.4pt,inner ysep=1.6pt,font=\scSmall,
                  text=scInk,minimum height=3.6mm},
  cofactor/.style={species,fill=scRefF,draw=scRef},
  catalyst/.style={species,fill=scHiF,draw=scHiE},
  rxn/.style={rectangle,fill=scRxF,draw=scRxE,line width=0.45pt,
              inner sep=1.2pt,font=\scSmall,text=scInk,
              minimum width=3.0mm,minimum height=3.0mm},
  place/.style={circle,fill=scSpF,draw=scSpE,line width=0.45pt,
                minimum size=4.0mm,inner sep=0pt},
  trans/.style={rectangle,fill=scRxE,draw=scRxE,minimum width=1.1mm,
                minimum height=3.6mm,inner sep=0pt},
  token/.style={circle,fill=scInk,minimum size=1.1mm,inner sep=0pt},
  % edges -----------------------------------------------------------------
  reac/.style={->,line width=0.5pt,draw=scReac,shorten >=0.8pt,shorten <=0.8pt},
  prod/.style={->,line width=0.5pt,draw=scProd,shorten >=0.8pt,shorten <=0.8pt},
  cat/.style={->,line width=0.5pt,draw=scHiE,shorten >=0.8pt,shorten <=0.8pt},
  flow/.style={-{Stealth[round,length=2.2mm,width=1.8mm]},line width=1.1pt,
               draw=scInkM!45},
  flab/.style={font=\scSmall,text=scInkM,inner sep=2pt},
  stoi/.style={font=\scSmall,text=scInkM,inner sep=0.8pt},
  % verdict marks: shape carries the meaning, colour only reinforces it
  vpass/.style={circle,fill=none,draw=none,text=black,
                font=\fontsize{6}{6}\selectfont\bfseries,
                minimum size=3.2mm,inner sep=0pt},
  vfail/.style={circle,fill=none,draw=none,text=black,
                font=\fontsize{6}{6}\selectfont\bfseries,
                minimum size=3.2mm,inner sep=0pt},
  callout/.style={font=\scSmall,text=scInkM,align=left},
}

% \scHead[<unused>]{<anchor coord>}{<letter>}{<title>}{<subtitle>}
% Plain "(a)" panel letter. Titles and subtitles are no longer printed: the
% caption names each panel. The arguments stay so the sources need no change.
\providecommand{\scHead}[5][scInk]{%
  \node[anchor=north west,inner sep=0pt,font=\scTitle] at #2
    {(\MakeLowercase{#3})};}

% --- shared pgfplots look (only if the figure loaded pgfplots) ------------
\ifdefined\pgfplotsset
  \pgfplotsset{compat=1.18}
  \pgfplotsset{
    synaxis/.style={
      scale only axis,
      tick align=outside, tick pos=left,
      axis line style={draw=scInkM!55,line width=0.4pt},
      % One grid, on the value axis only.  No minor grid: at these sizes it
      % turns the plot frame into a grey wash.
      ymajorgrids=true, xmajorgrids=false,
      major grid style={draw=scRuleL,line width=0.3pt},
      major tick length=1.6pt,
      label style={font=\scBody,color=scInk},
      tick label style={font=\scSmall,color=scInkM},
      legend style={font=\scSmall,draw=none,fill=none,inner sep=1pt,
                    row sep=0.4pt,legend columns=1},
      legend cell align=left,
      every axis plot/.append style={line width=0.7pt,mark size=1.2pt},
    },
  }
\fi

}
